\subsection{Researcher Positioning}

As a Black/African-American male scholar-practitioner, I must acknowledge that the framework this research presents (and much of the research itself) emerges from my specific positionality and cultural background. The development of Speculocultural Technopoiesis (ST) is deeply informed by Black epistemologies, cultural traditions, and theoretical perspectives (like Afrofuturism) that have shaped my own understanding of technology and creative practice. However, rather than presenting this as a limitation, I position this specific cultural grounding as a strength that offers valuable potential insights for addressing the shortcomings of many existing frameworks.

While the hope is that the core structure of the ST framework may be adaptable across different cultural contexts, my development and application of it is necessarily rooted in Black traditions and perspectives on sound and materiality. Thus, it is a \textbf{Black} \textit{Sonic} Speculocultural Technopoiesis. This approach resonates with Sylvia Wynter’s later articulation of ‘The Ceremony Found,’ which names the autopoetic act of re-envisioning the human beyond Western epistemic closure \parencite{Wynter_2015}. In this sense, ST engages in a kind of sonic "ceremony finding," a performative method that reveals the cultural specificity of universalized logics while opening space for alternative, speculative, and relational epistemologies.

This positioning reflects the framework's theoretical stance: \textit{that meaningful technological engagement must be culturally situated rather than universal in its abstraction}. The framework's structure is designed to be modular (in this context, meaning its components can be recombined and reconfigured to suit different cultural contexts while maintaining their functional integrity) and procedural (focusing on methodological processes and approaches rather than prescribing specific media formats or technological platforms). The hope is that this will allow practitioners from diverse cultural backgrounds to adapt its approach while maintaining deep engagement with their specific cultural contexts.

\subsection{Researcher Qualifications}

I approach this research as both practitioner and scholar. As an experimental media artist, designer, and musician whose creative practice informs a critical understanding of technology. Years of performance have taught me that precision (especially sonically) is always contextual to the craft, shaped by the body and environment. This embodied sensibility extends into the digital: each algorithm, interface, or system becomes an instrument that can be re-tuned through creative modification. In this sense, my practice turns technological design into a form of musicianship, where making, listening, and reworking are inseparable.

My academic training spans dual Bachelor’s degrees in Media Arts and Art Studio with concentrations in illustration, animation, and game design, and a Master of Art and Design in Experimental Media Arts focused on sonic interaction and listening as interaction. Graduate minors in Digital Humanities and Computer Science have provided interdisciplinary grounding in systems thinking, ludic design, and interface development. Technically, my expertise includes real-time engines (Unity and Unreal), node-based procedural tools (Houdini, Max/MSP, TouchDesigner), and programming for interactive audio-visual systems. These are the very tools ST seeks to unsettle and reimagine through cultural intervention.

Across my practice, research, and teaching, I work to connect embodied Black sonic epistemologies with computational logic. My projects and presentations at Ars Electronica (2023), the Atlantica Symposium on African Regroundings (2024), and the Media, Arts, and Design Conference (2024) reflect an ongoing commitment to research-creation as a mode of inquiry. This dual positioning as both artist and researcher enables me to integrate autoethnographic cultural knowledge with systematic technical rigor. Together, these capacities ground my ability to develop ST as both a theoretical framework and a practical methodology for culturally responsive, embodied design.